
\documentclass[12pt]{article}
\usepackage[margin=0.5in]{geometry} % Small margins for more text per page
\usepackage{fancyhdr}             % For custom headers/footers (page numbers)
\usepackage{listings}             % For formatting code
\usepackage{xcolor}               % For color in listings
\usepackage{titling}              % For title formatting

% Setup page numbering (fancyhdr prints page numbers by default)
\pagestyle{fancy}
\fancyhf{}
\rfoot{\thepage}

% Configure listings style for Python code
\lstset{
    language=Python,
    basicstyle=\ttfamily\small,
    breaklines=true,
    frame=single,
    numbers=left,
    numberstyle=\tiny\color{gray},
    keywordstyle=\color{blue},
    commentstyle=\color{green},
    stringstyle=\color{red},
    showstringspaces=false
}

\title{Retrieval Augmented Generation for Biomedical Database with AI Language Model}
\author{} % No author
\date{}   % No date

\begin{document}

% Cover page with title and description; suppress page number here if desired
\maketitle
\thispagestyle{empty}

\vspace{1cm}
\begin{center}
    \textbf{Description:} \\
    This script combines three approaches for Retrieval Augmented Generation (RAG) using multiple AI models (MedCPT-based encoders, llama\_index with OpenAI embeddings, and direct LLM query via llama\_index).
\end{center}

\newpage
\pagestyle{fancy}  % Reinstate page numbers after the cover page

% Insert the code file using the listings package
\lstinputlisting{rag-llm.py}

\end{document}
